\chapter[Embodied AI: The Loop Closes on the Factory Floor]{Embodied AI: The Loop Closes on the Factory Floor}
\chaptermark{Embodied AI}
\label{ch:embodied}

This chapter exists because its absence would be the most consequential omission the book could make. The thesis so far runs in one direction: China's manufacturing victories funded and shaped its digital-AI strategy. Embodied AI runs the loop the other way---model capability flowing back into factories, logistics, vehicles, laboratories, and physical services---and it is the domain where China's three structural advantages (manufacturing scale, deployment density, component supply chains) interact most directly with its open-model strategy. If the convergence thesis is correct anywhere with mechanical inevitability, it should be here, where the contested layer is not a benchmark but a bill of materials. And yet the chapter's conclusion is deliberately not a Chinese victory lap: the evidence as of this edition's date supports Chinese dominance of the \emph{hardware and deployment} layers and genuine uncertainty about whether hardware and deployment are the binding constraint at all. The question that decides the domain---is model capability or physical-system integration the harder bottleneck?---is precisely the question the rest of this book asks about digital AI, transposed into metal.

\section{The embodied stack, separated}
\label{sec:embodied-stack}

``Humanoid robot'' names a product, not an architecture, and the strategic analysis requires the stack:
\begin{align*}
\text{foundation model} \rightarrow \text{world model} \rightarrow{}& \text{task policy} \rightarrow \text{motion planning}\\
\rightarrow{}& \text{real-time control} \rightarrow \text{actuators and sensors},
\end{align*}
with a deterministic safety layer beside the whole chain and a fleet-learning loop above it. The layers have different physics, different competitive structures, and different nationalities. Language and vision reasoning (the foundation layer) is the commodity this book has spent four chapters on---and the embodied field increasingly consumes it in exactly the open form Part~II describes. World models---learned predictors of physical dynamics used for planning and for generating synthetic training experience---are a frontier-lab specialty (NVIDIA's Cosmos family, DeepMind's Genie line~\cite{cosmos,genie3}). Task policies are where the field's distinctive artifact lives: vision-language-action (VLA) models such as Physical Intelligence's $\pi$-class systems and Figure's Helix, typically a slow deliberative model (7--9\,Hz) coupled to a fast visuomotor policy (200\,Hz)~\cite{helix,pi-models}. Motion planning and real-time control are classical robotics, sixty years deep, dominated by Japanese and European engineering practice. Actuators and sensors are manufactured objects---which is to say, after Part~I, presumptively Chinese objects, a presumption the bill of materials below confirms.

The stack decomposition does the same analytical work the die-area argument did in Chapter~\ref{ch:compute}: it converts a category question (``who is winning humanoids?'') into a layer question, and the layer answers differ. The United States leads the top two layers. China leads the bottom two and the deployment environment. The middle---the task-policy layer where model meets metal---is genuinely contested, and it is the layer where the physical-data flywheel of the next section decides the outcome.

\section{The physical-data flywheel, and who is building it}
\label{sec:physical-flywheel}

Model the loop explicitly, because it is the third flywheel of this book and the one with the strangest property:
\begin{align*}
\text{deployment} \rightarrow \text{real interaction data} \rightarrow{}& \text{failure cases} \rightarrow \text{post-training}\\
\rightarrow{}& \text{better policy} \rightarrow \text{more deployment}.
\end{align*}
Web-scale text pretraining had the property that its input---humanity's accumulated writing---existed before the models did, and was free. Physical-interaction data has the opposite property: it does not exist until robots are deployed, it is generated in proportion to deployed fleet size and operating hours, and it cannot be scraped, only \emph{produced}. That single property is why the embodied contest may be structurally different from the LLM contest: the scarce input is not compute or algorithms but \emph{deployment surface}, and deployment surface is a function of manufacturing cost, component supply, factory density, and state coordination---the exact inventory of Chinese advantages from Part~I.

The Chinese state understood this early and built the loop's infrastructure as public works. By late 2025 more than forty state-backed robot data-collection centers had been announced (roughly two dozen operational), staffed by human teleoperators in VR headsets and exoskeletons demonstrating assembly, household, and elder-care tasks for hours a day---industrialized imitation-learning data production [V]~\cite{data-farms}. Shanghai's heterogeneous training facility opened in January 2025 with capacity for a hundred robots simultaneously and a target of a thousand by 2027 collecting ten million entries~\cite{kylin}; Beijing's Shijingshan camp runs sixteen scenario sets from car-assembly lines to smart homes; provincial centers have become a demand channel in their own right, with UBTech booking RMB\,566M of robot sales to three of them and China Mobile placing nine-figure orders with Unitree and AgiBot~\cite{data-farms}. On the open side, AgiBot's World Colosseo dataset publishes a million-plus real-robot trajectories from a hundred robots across a hundred scenarios---the embodied analogue of the open-weight release, and aimed at the same strategic effect~\cite{agibot-world}. The Western equivalents are real but differently shaped: Open X-Embodiment pools academic lab data across 33 institutions~\cite{open-x}; Figure trained Helix on roughly five hundred hours of curated teleoperation~\cite{helix}; Agility has accumulated 65{,}000+ hours of genuinely productive deployment~\cite{agility-spac}; Tesla possesses, in principle, the largest fleet-learning apparatus on earth and has yet to ship a production humanoid to use it on~\cite{optimus-zero}.

The honest current reading: China is industrializing the \emph{production} of physical-interaction data; the United States leads the \emph{consumption} of it (the VLA models that turn it into capability). Neither half works without the other, which is why this flywheel is the one whose ownership is least settled---and why it cannot be rented: a laboratory without manufacturing access cannot buy deployment surface the way it can buy tokens, and a manufacturer without frontier post-training cannot convert demonstrations into generalization. The strategic symmetry with Chapter~\ref{ch:theory}'s two thefts is exact, and Section~\ref{sec:embodied-positions} prices it.

\section{The industrial bill of materials}
\label{sec:embodied-bom}

The component layer is where Part~I's precedents repeat with the least resistance, because it \emph{is} Part~I's industries wearing new part numbers. Table~\ref{tab:embodied-bom} scores the major component classes.

\begin{table}[htbp]
  \centering
  \scriptsize
  \caption{The embodied bill of materials, mid-2026 [A unless noted]. ``Position'' names the current center of gravity, not a monopoly.}
  \label{tab:embodied-bom}
  \setlength{\tabcolsep}{3.5pt}
  \begin{tabular}{p{30mm}p{52mm}p{54mm}}
    \toprule
    Component & Position & Evidence and trajectory \\
    \midrule
    Strain-wave (harmonic) reducers & Japan incumbent (Harmonic Drive); China scaling fast & Leaderdrive is China's largest maker, lifted by humanoid demand; Zhongda Leader's founders became billionaires on joint demand; Unitree's price war is already squeezing the domestic reducer chain~\cite{reducers} \\
    \addlinespace
    Planetary roller screws & China dominant in new capacity & $\sim$33\% of humanoid BoM by one bank estimate; unit cost \$1,350--2,700; Shanghai Beite's \$260M dedicated plant; thin Western capacity~\cite{roller-screws} \\
    \addlinespace
    Actuator assemblies, servo \& frameless motors & China (three parallel chains: Tesla-, Huawei-, Unitree-anchored) & reported \$685M Optimus actuator order to Sanhua; Tuopu building dedicated drive-system capacity~\cite{tesla-suppliers} \\
    \addlinespace
    Lidar and depth sensing & China decisively & Hesai robotics shipments $+$744\% y/y (49k units, Q2 2025); RoboSense $+$632\%; both pivoting from automotive~\cite{lidar-robotics} \\
    \addlinespace
    Force/torque, tactile, dexterous hands & Contested; APAC-centered manufacturing & dozens of Chinese hand makers; Japanese and German precision sensing incumbents \\
    \addlinespace
    Batteries & China (Chapter~\ref{ch:batteries}) & CATL both supplies packs and deploys humanoids on its own pack lines---the loop in one firm~\cite{catl-galbot} \\
    \addlinespace
    Edge compute & US platform (Jetson Thor, \$3,499, 2{,}070 FP4 TFLOPS); domestic chain behind & Huawei/Cambricon alternatives exist inside the Chinese chains~\cite{jetson-thor} \\
    \addlinespace
    Safety systems \& certification & Europe and Japan & ISO 10218 (2025 ed.), ISO 25785-1 for dynamically stable mobile robots drafting toward $\sim$2027~\cite{iso-humanoid} \\
    \addlinespace
    Simulation \& digital twins & US (Isaac/Cosmos); China open-sourcing fast (Genie Sim) & \cite{cosmos,agibot-world} \\
    \bottomrule
  \end{tabular}
\end{table}

Two cost-curve observations convert the table into strategy. First, the slope: Goldman Sachs found humanoid unit manufacturing cost falling from \$50--250k to \$30--150k in a single year---a $\sim$40\% decline where 15--20\% was forecast---and Morgan Stanley finds localized Chinese supply chains running $\sim$20\% cheaper than foreign ones, with average system prices trending from $\sim$\$46k toward $\sim$\$21k [A]~\cite{goldman-humanoid,ms-humanoid}. Unitree's own ASP fell from RMB\,593k (2023) to RMB\,166k (2025), with an entry model at RMB\,29{,}900---the solar learning curve, replayed on legs, with the familiar accompaniment of collapsing gross margins (87.7\%\,$\rightarrow$\,63.2\%) that Part~I teaches us to read as the cull beginning~\cite{unitree-ipo}. Second, the dependency: supply-chain reporting places on the order of 70\% of Tesla Optimus's supply chain in China, with estimates that excluding Chinese parts would triple its cost [A]---and China demonstrated the leverage directly when rare-earth export licensing touched Optimus's actuator magnets in 2025~\cite{optimus-china}. The American humanoid, like the American solar panel of 2010, is being designed atop the adversary's supply chain. This is playbook move~M5 (supply-chain capture) achieved \emph{before} the product category even exists at scale---a first even for China.

\section{The deployment scoreboard, and what it does not show}
\label{sec:embodied-scoreboard}

The shipment table reads like a rout. Global humanoid shipments exceeded 13{,}000 units in 2025, roughly 85\% from Chinese firms [A]~\cite{fortune-humanoid}. AgiBot shipped its 10{,}000th cumulative unit in March 2026---its second five thousand in roughly three months~\cite{agibot-10k}. Unitree sold 5{,}500 humanoids in 2025 and, almost uniquely in the industry worldwide, did so \emph{profitably} ($\sim$\$250M revenue, $\sim$\$41M profit), reaching its STAR-market IPO approval in June 2026~\cite{unitree-ipo}. UBTech put its thousandth Walker S2 off the line with cumulative orders past RMB\,800M and multi-robot deployments at Zeekr, BYD, Foxconn, and Audi-FAW plants~\cite{ubtech}. CATL began scale deployment of Galbot systems on its own battery lines~\cite{catl-galbot}; BYD's four-year internal program targets $\sim$20{,}000 robots in its own factories in 2026 [P]~\cite{byd-robots}; XPeng targets mass production of its Iron humanoid by end-2026 and a million units by 2030 [R]~\cite{xpeng-iron}. Around them, the ecosystem markers of every Part~I campaign: 150-plus companies, embodied AI named in the 2025 Government Work Report and the 15th Five-Year Plan, an \$8.2B national fund allocation, Beijing's \$14.3B fifteen-year robotics fund, provincial subsidy programs, and municipal action plans with unit and revenue targets~\cite{miit-humanoid,embodied-policy}.

The American column is shorter and stranger. Tesla---the company whose CEO calls Optimus ``the biggest product ever''---had zero verified production units as of July 2026, deferred its Gen-3 reveal twice, and told its own earnings call that existing units are ``primarily for learning, not productive tasks''~\cite{optimus-zero}. Figure, valued at \$39B, has delivered 350-plus units of its third-generation robot and produced the single most important capability datapoint in the Western column: a livestreamed eight-hour \emph{autonomous} logistics shift at human-comparable package-sorting rates~\cite{helix,figure-c}. Agility's Digit has 65{,}000-plus hours of real deployed operation and a \$300M+ multi-year order book, going public via SPAC at \$2.5B~\cite{agility-spac}. Boston Dynamics' electric Atlas has its entire 2026 production run committed and a Hyundai plant designed for 30{,}000 units a year---gated, instructively, by a Korean labor union blocking deployment absent an agreement~\cite{atlas}. Behind them stand the model shops: Physical Intelligence at a reported $\sim$\$5.6B raising toward \$11B [A], Skild at a reported $\sim$\$14B [A], and NVIDIA supplying the compute, the world models, and an open reference humanoid to the entire field~\cite{pi-models,jetson-thor}.

Now the deflationary reading, which this book owes the reader after four chapters of proof-boundary discipline. \emph{Shipments are not labor.} The majority of Chinese humanoid units go to demonstrations, exhibitions, education, and---in a perfect circularity---to the state-funded data-collection centers whose purchases are themselves industrial policy; independent assessments describe most deployed units as ``performative rather than functional,'' fragile in operation and dependent on structured environments~\cite{fortune-humanoid}. Teleoperation shadows the entire field's marketing, on both sides of the Pacific: Tesla's bartending robots were remotely operated, the 1X NEO home robot ships with human teleoperators explicitly in the loop, and the viral demonstrations that anchor valuations routinely turn out to have a human in them~\cite{teleop}. The comparison that matters is therefore not units shipped but \emph{intervention-free productive hours}, and on that metric the public record inverts the shipment table: the best-documented sustained autonomous performances (Figure's eight-hour shift; Agility's cumulative deployment record) are American~\cite{helix,agility-spac}. The scoreboard, graded honestly: China leads manufacturing scale, cost, components, and deployment surface [V]; the United States leads demonstrated autonomy per deployed unit [V]; and the metric that would settle the contest---useful work per robot-hour at fleet scale---is published by nobody.

\section{Metrics that matter: the embodied proof suite}
\label{sec:embodied-metrics}

As with the LX2 in Chapter~\ref{ch:compute}, the remedy for a contested frontier is a stated measurement program. The quantities that would decide the embodied contest, none of which any major player currently publishes systematically: intervention-free operating hours; task-success rate on standardized task sets; cycle time against human baseline; uptime; mean time between safety stops; human-supervision minutes per robot-hour (the teleoperation ratio, the single most information-dense number in the field); cost per successful physical task; energy per task; damage and near-miss rates; skill-transfer performance across sites (train in one factory, deploy in another---the sim-to-real and site-to-site generalization that separates a product from a demo); and useful training data generated per deployed unit, which is the flywheel's own gauge. A vendor that publishes this table for a hundred-robot deployment over six months will do for embodied AI what MLPerf did for accelerators; until one does, every shipment figure in the previous section should be read as capacity, not output. The forecast register (Appendix~\ref{app:dashboard}) now carries the corresponding entries.

\section{Comparative national positions}
\label{sec:embodied-positions}

The five-way comparison, stated without a predetermined winner. \textbf{China} holds manufacturing scale, the component base, deployment density (54\% of world industrial-robot installations, 57\% domestic-brand share of its own market, the largest operational stock ever recorded~\cite{ifr}), state-industrialized data collection, and the cost curve; its exposure is the familiar pair---capability of the top model layers, and a domestic price war already compressing margins before the product works. \textbf{The United States} holds the frontier models, the world models, the VLA research lead, the capital (three embodied startups above \$10B), and the strongest demonstrated autonomy; its exposure is a supply chain it does not own and a deployment environment of high labor cost but low robot density. \textbf{Japan} holds the precision-component incumbency (strain-wave gearing above all), the industrial-robot giants, production engineering as a craft, and---via Chapter~\ref{ch:strategies}'s logic---a natural niche as the trusted supplier of the components both blocs need; its exposure is having invented the category (ASIMO) and ceded the product. \textbf{Korea} holds the Hyundai--Boston Dynamics vertical (a 30k-unit plant attached to a captive automaker) plus its own robotics base. \textbf{Europe} holds the safety-certification tradition and industrial-integration expertise, and---true to form---is currently converting them into standards rather than products.

The decisive uncertainty is which bottleneck binds. If \emph{model capability} is the constraint---if generalist manipulation at production reliability is still several breakthroughs away---then the American position is stronger than the shipment tables suggest, Chinese fleets amortize into demos, and the data farms are collecting demonstrations for policies that cannot yet use them. If \emph{physical-system integration} is the constraint---if current VLAs are good enough that reliability, cost, and iteration speed at fleet scale decide---then China's position compounds exactly as it did in EVs, and the American labs are optimizing the layer that is about to commoditize. The evidence is genuinely mixed: the eight-hour autonomous shift argues capability has arrived for narrow logistics; the near-universal teleoperation dependence argues it has not for anything else. This book's judgment [S], registered as a forecast: the integration reading wins in \emph{structured} environments (factories, warehouses) within the decade---China's home terrain---while unstructured environments (homes, care, field work) remain capability-bound well past 2030, leaving the largest advertised markets unclaimed by anyone.

\section{Falsifiers}
\label{sec:embodied-falsifiers}

Consistent with the book's method, the chapter states what would prove its cautious-Chinese-advantage reading wrong in either direction. \emph{Against the embodied thesis entirely}: teleoperation ratios that do not decline across 2026--2029; intervention-free hours that stagnate at demo scale; simulated training that persistently fails to transfer; safety certification (ISO 25785-class) blocking broad deployment the way certification blocked the C919; unit economics that stay worse than human labor in every unstructured setting; hardware reliability, not intelligence, remaining the binding constraint---humanoids as the Concorde of automation. \emph{Against the Chinese-advantage reading specifically}: a Western lab achieving generalist manipulation that renders demonstration-data volume irrelevant (the capability escape); US/allied reshoring of the actuator and reducer chains at competitive cost; the Chinese price war culling the sector before any firm reaches sustainable unit economics---involution consuming the industry in its infancy, the failure mode Section~\ref{sec:china-advice} warns Beijing about directly; or the state data-farm model producing large volumes of low-diversity demonstrations that generalist policies turn out not to need. Each has an observable signature, and the quarterly dashboard now watches for all of them.

The chapter's closing symmetry deserves a sentence, because it completes the book's architecture. Part~I showed physical manufacturing funding China's route to digital AI. This chapter shows digital AI attempting to conquer physical manufacturing. If both directions close, the two flywheels of Chapter~\ref{ch:theory} fuse into a single self-reinforcing industrial system---AI improving the factories that build the machines that generate the data that improves the AI---and the country that owns that closed loop owns something qualitatively beyond anything Part~I described. That, and not any single robot, is the stake on the factory floor.
